\section{Conclusion}
\label{sec:conclusion}

The Article has argued that delegated emergency authority across five doctrinally disparate regimes shares
a structural feature whose absence explains a recurring failure mode.
The feature is the requirement that an authority, at the moment of its exercise, exhibit a typed witness
for the path back to the prior state.
The feature is not currently a part of any of the five statutory or constitutional regimes examined.
Its absence is, on the Article's framing, the structural explanation for the ratcheting pattern documented
in the historical record of each regime.

The structural correction the Article proposes is a formal grammar drawn from a Lean-attestable substrate
developed in a companion technical paper.
The grammar requires four components for any bounded executive action:
the substantive action,
a time-to-live,
a typed rollback witness, and
a quorum predicate.
The grammar makes the unconstructive case unconstructible:
an action that cannot exhibit a typed rollback witness at the moment of admission cannot be admitted.
The discipline operates at the construction step, not at a later review step.

The Article has not argued that any particular delegated emergency authority is desirable or undesirable on the merits.
It has argued that, structurally, an emergency authority enacted without the four components is a defective grant of authority,
and that the defect has substantive consequences that the historical record documents.

Three threads of further work are visible.
The first is the engagement of constitutional-law scholarship in detail.
The Article's case studies are written by a non-lawyer working from secondary literature.
A constitutional-law scholar would write each case study differently, and the variation would be informative.
The second is the engagement of the structural correction with the doctrinal toolkit of sunset clauses,
periodic reauthorization, and supermajority defaults.
The Article has argued that the typed-rollback discipline is a genuine addition to that toolkit rather than a relabeling,
but the comparative analysis of when the direct address is preferable to the indirect addresses is open.
The third is the application of the grammar to delegated emergency authority in regimes the Article has not
examined.
Lincoln's habeas suspensions, Article~16 of the French Fifth Republic, the Italian decree-law tradition,
the Indian President's Rule, and the Brazilian state of defense regime are all instances of the pattern.
Each would test the grammar's robustness against doctrinal variation the Article has not engaged.

The Article's central contribution is the identification of the pattern,
the naming of the structural correction,
and the demonstration that the correction is buildable.
The substantive content of any application of the grammar is a matter of democratic political determination
that the grammar deliberately leaves to the political process.
What the grammar offers is a structural skeleton within which that determination is made.
The skeleton, on the Article's argument, is what the doctrinal record of delegated emergency authority has been missing.
